\documentclass{article}
\usepackage{amsmath, amssymb, amsthm}
\usepackage{hyperref}
\usepackage{geometry}
\geometry{margin=1in}

\title{Erd\H{o}s Problem \#1122}
\date{Last updated: 2025-12-30}
\author{Source: erdosproblems.com}

\begin{document}
\maketitle

\noindent\textbf{Status:} OPEN \quad \textbf{No prize}

\noindent\textbf{Tags:} number theory

\medskip

\noindent\textbf{Problem Statement:}

Let $f:\mathbb{N}\to \mathbb{R}$ be an additive function (i.e. $f(ab)=f(a)+f(b)$ whenever $(a,b)=1$). Let\[A=\{ n \geq 1: f(n+1)&#60; f(n)\}.\]If $\lvert A\cap [1,X]\rvert =o(X)$ then must $f(n)=c\log n$ for some $c\in \mathbb{R}$?




Erd\H{o}s proved that $f(n)=c\log n$ for some $c\in\mathbb{R}$ if $A$ is empty, or if $f(n+1)-f(n)=o(1)$. Partial progress was made by Mangerel \cite{Ma22}, who proved that this is true if\[\lvert A\cap [1,X]\rvert \ll \frac{X}{(\log X)^{2+c}}\]for some $c&gt;0$, and if $f(p)$ does not have very large values (in a certain technical sense). See also [491] .




References

[Ma22] Mangerel, Alexander P., Additive functions in short intervals, gaps and a conjecture of {E}rd\H{o}s . Ramanujan J. (2022), 1023--1090.


\bigskip
\noindent\small{Source: \url{https://www.erdosproblems.com/1122}}
\end{document}
